% Generated from tex/chapters/ch04/06_構築の管理.tex for the SWEBOK structure tree
\subsubsection{構築測定}

構築では、多くの活動や成果物を測定できる。開発したコード量、変更したコード量、再利用したコード量、削除したコード量、コード複雑度、コードレビューで見つかった欠陥数、欠陥発見率、欠陥修正率、作業工数、予定との差分などである。

測定の目的は、管理と改善である。数値を集めるだけでは意味がない。たとえば、複雑度が高い箇所を重点的にレビューする、欠陥修正率から計画を見直す、再利用量から開発効率を評価する、といった使い方が必要である。
